\documentclass[11pt]{article}
\usepackage[letterpaper,margin=1in]{geometry}
\usepackage[T1]{fontenc}
\usepackage[utf8]{inputenc}
\usepackage{lmodern,microtype}
\usepackage{amsmath,amssymb,mathtools,bm,slashed}
\usepackage{xcolor,enumitem,booktabs,array,needspace}
\usepackage[most]{tcolorbox}
\usepackage{fancyhdr,titlesec,parskip,hyperref}
\definecolor{qftblue}{HTML}{245A8D}
\definecolor{qftgreen}{HTML}{2F6B4F}
\definecolor{qftgray}{HTML}{F4F6F8}
\hypersetup{colorlinks=true,linkcolor=qftblue,urlcolor=qftblue,pdfauthor={Nima Lashkari}}
\pagestyle{fancy}\fancyhf{}
\lhead{PHYS 662: Quantum Field Theory I}\cfoot{\thepage}
\setlength{\headheight}{14pt}
\titleformat{\section}{\Large\bfseries\color{qftblue}}{\thesection.}{0.6em}{}
\titleformat{\subsection}{\large\bfseries\color{qftgreen}}{\thesubsection}{0.6em}{}
\setlist[enumerate]{leftmargin=2em,itemsep=0.25em,topsep=0.35em}
\setlist[itemize]{leftmargin=2em,itemsep=0.2em,topsep=0.35em}
\newcommand{\dd}{\mathrm d}
\newcommand{\ii}{\mathrm i}
\newcommand{\eps}{\varepsilon}
\newcommand{\tr}{\operatorname{tr}}
\newcommand{\Tr}{\operatorname{Tr}}
\newcommand{\diag}{\operatorname{diag}}
\newcommand{\Lag}{\mathcal L}
\newcommand{\D}{\mathcal D}
\newcommand{\vac}{\lvert0\rangle}
\newcommand{\vbra}{\langle0\rvert}
\newcommand{\Tord}{\mathrm T}
\newcommand{\ket}[1]{\lvert#1\rangle}
\newcommand{\bra}[1]{\langle#1\rvert}
\newtcolorbox{keyidea}[1][]{enhanced,breakable,colback=qftgray,colframe=qftblue,
 boxrule=0.7pt,arc=1.5mm,title={#1},fonttitle=\bfseries,coltitle=white}
\newtcolorbox{remarkbox}[1][]{enhanced,breakable,colback=white,colframe=qftgreen,
 boxrule=0.6pt,arc=1.5mm,title={#1},fonttitle=\bfseries,coltitle=white}
\newcommand{\intp}{\int\!\frac{\dd^4p}{(2\pi)^4}}
\newcommand{\MSbar}{\overline{\mathrm{MS}}}
\newcommand{\module}[1]{\clearpage\section{#1}}
\newcommand{\lecturefront}[4]{%
\rhead{Lecture #1}\hypersetup{pdftitle={PHYS 662 Lecture #1: #2}}
\begin{document}\pagenumbering{gobble}
\begin{titlepage}\centering\vspace*{1.05in}
{\scshape\Large PHYS 662 --- Quantum Field Theory I\par}
\vspace{0.35in}{\Huge\bfseries\color{qftblue}Lecture #1\par}
\vspace{0.2in}{\LARGE\bfseries #2\par}
\vspace{0.2in}{\large\color{qftgreen}#3\par}
\vspace{0.6in}{\large Nima Lashkari\par}{\large Fall 2026\par}
\vfill\begin{minipage}{0.86\textwidth}\small #4\par\medskip
Course page: \url{https://nimalashkari.com/quantum-field-theory-2026-part-i/}
\end{minipage}\vspace{0.5in}\end{titlepage}
\pagenumbering{roman}\setcounter{tocdepth}{1}\tableofcontents
\vfill\begin{remarkbox}[Place in the course]
Lectures 19--20 established free scalar and fermion path integrals.
Lecture 21 turns the scalar source expansion into Feynman rules.
Lecture 22 evaluates the first scalar loops and introduces renormalization.
Problem Set 7 connects the diagram counting to these calculations.
\end{remarkbox}\clearpage\pagenumbering{arabic}}
\lecturefront{21}{Interactions in the Path-Integral Formalism}
{From source derivatives to Feynman rules}
{We use the free scalar Gaussian as a calculational tool for an interacting
theory. The main goals are to derive diagram weights, distinguish
connected and one-particle-irreducible functions, and translate between
position space, momentum space, and scattering amplitudes.}

\section{An interacting theory built from the free Gaussian}
The scalar source functional of Lecture 19 already contains all free
correlation functions. An interaction changes their weights and connects
contractions that were previously independent. We work in the symmetric
vacuum of a massive real scalar, with \(m^2>0\), small positive \(\lambda\),
\(\hbar=c=1\), and \(\eta_{\mu\nu}=\diag(1,-1,-1,-1)\):
\begin{equation}
 S[\phi]=S_0[\phi]+S_{\rm int}[\phi],\qquad
 \Lag=\frac12(\partial\phi)^2-\frac12m^2\phi^2-\frac{\lambda}{4!}\phi^4,
 \quad S_{\rm int}=-\frac{\lambda}{4!}\int\dd^4x\,\phi^4.
 \label{eq:action}
\end{equation}
The minus sign in \(S_{\rm int}\) corresponds to a positive quartic
potential. The factor \(4!\) is a useful convention, not an extra
physical parameter. At this stage the theory is regulated; coincident
contractions will be defined and renormalized in Lecture 22.

Normalize the vacuum functional from the beginning:
\begin{equation}
 Z[J]=\frac{\int\D\phi\,e^{\ii S[\phi]+\ii\int J\phi}}
 {\int\D\phi\,e^{\ii S[\phi]}},\qquad Z[0]=1,\qquad
 G_n(x_1,\ldots,x_n)=\left.\frac{1}{\ii^n}
 \frac{\delta^n Z}{\delta J(x_1)\cdots\delta J(x_n)}\right|_{J=0}.
 \label{eq:Z}
\end{equation}
The \(\ii0\) prescription, or a small imaginary-time tilt at the
endpoints, selects the interacting vacuum \(\ket{\Omega}\).
Thus \(G_n=\langle\Omega|\Tord\phi(x_1)\cdots\phi(x_n)|\Omega\rangle\).
This is a vacuum correlator, not a thermal trace or a fixed-field
transition kernel. Those other boundary conditions give other Green
functions even for the same action.

Keep the propagator convention of the previous lectures:
\begin{equation}
 D_F(x-y)=\intp\frac{\ii e^{-\ii p\cdot(x-y)}}{p^2-m^2+\ii0},
 \qquad (\Box_x+m^2)D_F(x-y)=-\ii\delta^{(4)}(x-y),
 \label{eq:DF}
\end{equation}
\begin{equation}
 Z_0[J]=\exp\!\left[-\frac12\int\dd^4x\dd^4y\,
 J(x)D_F(x-y)J(y)\right].                         \label{eq:freeZ}
\end{equation}
\begin{keyidea}[One convention for every diagram]
A contraction is \(D_F\), including its numerator \(\ii\).
A quartic vertex will be \(-\ii\lambda\).
Do not attach another \(\ii\) to a line. Odd vacuum correlators vanish
because we expand about the unbroken \(\phi\mapsto-\phi\) vacuum.
\end{keyidea}
The denominator of \eqref{eq:Z} removes components with no external
insertions. It does not make every remaining diagram connected.

\module{Exact equations before the expansion}
Functional integration by parts gives a useful check on perturbation
theory. With a regulator in place and boundary terms absent,
\begin{equation}
 0=\int\D\phi\,\frac{\delta}{\delta\phi(x)}
 \left(F[\phi]e^{\ii S[\phi]}\right)
 \quad\Longrightarrow\quad
 \left\langle F\frac{\delta S}{\delta\phi(x)}\right\rangle
 =\ii\left\langle\frac{\delta F}{\delta\phi(x)}\right\rangle.
 \label{eq:ibp}
\end{equation}
Using a source exponential as \(F\), and
\(\delta S/\delta\phi=-(\Box+m^2)\phi-\lambda\phi^3/3!\), gives
\begin{equation}
 \left[(\Box_x+m^2)\frac{1}{\ii}\frac{\delta}{\delta J(x)}
 +\frac{\lambda}{3!}\left(\frac{1}{\ii}\frac{\delta}{\delta J(x)}\right)^3
 -J(x)\right]Z[J]=0.                              \label{eq:SD}
\end{equation}
This is the Schwinger--Dyson equation. At zero source it relates the
one-point and three-point functions. One additional source derivative
produces a contact term:
\begin{equation}
 (\Box_x+m^2)G_2(x,y)+\frac{\lambda}{3!}G_4(x,x,x,y)
 =-\ii\delta^{(4)}(x-y).                          \label{eq:SD2}
\end{equation}
The delta function is essential. Time ordering and coincident operator
insertions prevent us from using the classical equation inside a
correlator without contact terms.

For \(X=(x_1,\ldots,x_n)\), let \(X_{\widehat a}\) omit \(x_a\).
Repeated differentiation gives the hierarchy
\begin{align}
 (\Box_x+m^2)G_{n+1}(x,X)+\frac{\lambda}{3!}G_{n+3}(x,x,x,X)
 &=-\ii\sum_{a=1}^n\delta^{(4)}(x-x_a)G_{n-1}(X_{\widehat a}),
 \label{eq:hierarchy}\\
 G_{n+1}(x,X)&=\sum_{a=1}^nD_F(x-x_a)G_{n-1}(X_{\widehat a})\notag\\
 &\quad-\frac{\ii\lambda}{3!}\int\dd^4z\,
 D_F(x-z)G_{n+3}(z,z,z,X).                        \label{eq:integralSD}
\end{align}
The last line uses the Feynman inverse and vacuum boundary conditions.
For \(\lambda=0\), it recursively generates Wick's pairings, with
\(G_0=1\). For nonzero coupling it is an infinite coupled system, not
a closed equation for \(G_2\).

\begin{remarkbox}[A mean field is not the full quantum equation]
The identity \(\langle V'(\phi)\rangle\) is generally different from
\(V'(\langle\phi\rangle)\). Replacing the former by the latter discards
fluctuations and requires an approximation. The symmetric vacuum has
\(\langle\phi\rangle=0\); a general interacting one-point function is
not determined by the extrema of the classical potential alone.
\end{remarkbox}

\module{Perturbation theory and the cancellation of vacuum bubbles}
Every field in the interaction can be replaced by a source derivative.
Define the interaction differential operator
\[
 \widehat V=-\frac{\ii\lambda}{4!}\int\dd^4z\,
 \left(\frac{1}{\ii}\frac{\delta}{\delta J(z)}\right)^4.
\]
Then the normalized expansion is
\begin{equation}
 Z[J]=\frac{e^{\widehat V}Z_0[J]}
 {\left.e^{\widehat V}Z_0[J]\right|_{J=0}}.         \label{eq:operatorZ}
\end{equation}
Expanding the exponential organizes the answer by the number of
vertices. In general this is an asymptotic expansion; we do not assume
convergence of the infinite series.

A first-order calculation shows what the normalization does. Write
\begin{equation}
 f(z)=\int\dd^4x\,D_F(z-x)J(x),\qquad
 \frac{1}{\ii Z_0}\frac{\delta Z_0}{\delta J(z)}=\ii f(z).
\end{equation}
The Gaussian with source has mean \(\ii f\) and connected covariance
\(D_F\). Four derivatives therefore give
\begin{equation}
 \frac{1}{Z_0}\left(\frac{1}{\ii}\frac{\delta}{\delta J(z)}\right)^4Z_0
 =3D_F(0)^2-6D_F(0)f(z)^2+f(z)^4.                \label{eq:polynomial}
\end{equation}
There are three ways to pair four fields, six ways to choose the two
fields replaced by their mean, and one term with four means. The
relative minus sign follows from \((\ii f)^2=-f^2\).

The term \(3D_F(0)^2\) is independent of \(J\), so the denominator of
\eqref{eq:operatorZ} cancels it:
\begin{equation}
 Z[J]=Z_0[J]\left\{1-\frac{\ii\lambda}{4!}\int\dd^4z\,
 [f(z)^4-6D_F(0)f(z)^2]+O(\lambda^2)\right\}.    \label{eq:Zfirst}
\end{equation}
Two source derivatives now yield
\begin{equation}
 G_2(x,y)=D_F(x-y)-\frac{\ii\lambda}{2}D_F(0)
 \int\dd^4z\,D_F(x-z)D_F(z-y)+O(\lambda^2).
 \label{eq:tadpoleposition}
\end{equation}
The factor \(1/2\) is the tadpole's diagram weight. Applying
\(\Box_x+m^2\) to this expression reproduces \eqref{eq:SD2} with
\(G_4^{(0)}(x,x,x,y)=3D_F(0)D_F(x-y)\), independently checking the sign.

At higher orders, components with no external points exponentiate
and factor out of the numerator. The denominator cancels that entire
factor. Disconnected products attached to external points survive:
for example the three free pairings of a four-point function.
Vacuum-energy calculations require the unnormalized functional;
normalization deliberately removes that information here.

\module{Wick contractions and symmetry factors}
At order \(V\), the interaction expansion contains
\(1/[V!(4!)^V]\). Wick contraction produces many identical integrals.
For a graph with fixed, labeled external points, its coefficient is
\begin{equation}
 \frac{N_{\rm contractions}}{V!(4!)^V}\equiv\frac{1}{S}.
 \label{eq:symmetry}
\end{equation}
Here \(S\) is the symmetry factor. It counts relabelings of internal
vertices and line ends that leave the graph unchanged while keeping
external labels fixed. A reflection that exchanges two distinct
external particles is not automatically such a symmetry.

\begin{center}
\begin{tikzpicture}[scale=0.95,thick,every node/.style={font=\small}]
 \draw (0,0) circle (0.4);\draw (0,0.8) circle (0.4);
 \fill (0,0.4) circle (2pt);\node at (0,-0.75){vacuum: \(S=8\)};
 \draw (2,0)--(4,0);\draw (3,0.4) circle (0.4);
 \fill (3,0) circle (2pt);\node[left] at (2,0){\(x\)};
 \node[right] at (4,0){\(y\)};\node at (3,-0.75){tadpole: \(S=2\)};
 \begin{scope}[xshift=1cm]
 \draw (5.2,-0.25)--(6,0.4)--(6.8,-0.25);
 \draw (5.2,1.05)--(6,0.4)--(6.8,1.05);
 \fill (6,0.4) circle (2pt);\node at (6,-0.75){four-point tree: \(S=1\)};
 \end{scope}
\end{tikzpicture}
\end{center}
For the one-vertex vacuum graph there are three pairings, so
\(3/4!=1/8\). For the tadpole, attach the first external field to one
of four vertex fields and the second to one of the remaining three.
The last two contract with each other: \(4\cdot3/4!=1/2\).
For a four-point tree, all four external fields attach to the vertex:
\(4!/4!=1\).

A less immediate example is a four-point bubble with the partition
\((12|34)\). Two lines join its two vertices. Either integrated vertex
can receive \(x_1,x_2\), and each receives its two external fields in
\(4\cdot3\) ways. The remaining ends join in \(2!\) ways. Consequently
\begin{equation}
 \frac{2(4\cdot3)^2\,2!}{2!(4!)^2}=\frac12.
 \label{eq:bubblecount}
\end{equation}
There are three partitions, \((12|34),(13|24),(14|23)\). Each has
weight \(1/2\); this is separate from summing the three channels.

\begin{keyidea}[Position-space rules for the real scalar]
Draw the desired graphs with labeled external points. Put \(D_F(x-y)\)
on every line, \(-\ii\lambda\) at each quartic vertex, and integrate
each vertex over \(\dd^4z\). Multiply by \(1/S\). For normalized
correlators exclude vacuum components; for connected correlators keep
only connected graphs. Include each distinct external assignment once.
\end{keyidea}
These rules abbreviate a derivation from the Gaussian. When a diagram
weight is uncertain, counting contractions is more reliable than
judging the apparent symmetry of a particular drawing.

\module{Connected functions and irreducible interaction kernels}
Define \(W[J]=-\ii\log Z[J]\), as in Lecture 19. The exponential formula
for disconnected graphs gives
\begin{equation}
 \log Z[J]=\sum_{n\geq1}\frac{\ii^n}{n!}\int J(x_1)\cdots J(x_n)
 G_{n,c}(x_1,\ldots,x_n),\qquad
 \left.\delta_J^nW\right|_0=\ii^{n-1}G_{n,c}.     \label{eq:connected}
\end{equation}
Spacetime integrations are implicit in the first expression. In the
symmetric vacuum \(G_{2,c}=G_2\), while
\begin{equation}
 G_4(1,2,3,4)=G_{4,c}(1,2,3,4)
 +G_2(1,2)G_2(3,4)+G_2(1,3)G_2(2,4)+G_2(1,4)G_2(2,3).
 \label{eq:G4decompose}
\end{equation}
The two-point functions in this identity are the full interacting ones.
At leading order the new connected part is
\begin{equation}
 G_{4,c}^{(1)}(x_1,x_2,x_3,x_4)
 =-\ii\lambda\int\dd^4z\,\prod_{a=1}^4D_F(x_a-z).
 \label{eq:tree4}
\end{equation}
A connected graph is \emph{one-particle irreducible} (1PI) if cutting any
one internal line cannot disconnect it. A self-energy insertion on an
external leg of \eqref{eq:tree4} is connected but one-particle reducible.
The two lines in a bubble provide two paths between its vertices;
cutting just one leaves it connected.

\begin{center}
\begin{tikzpicture}[thick,every node/.style={font=\small}]
 \draw (-0.8,-0.6)--(0,0)--(-0.8,0.6);
 \draw (2.3,-0.6)--(1.5,0)--(2.3,0.6);
 \draw (0,0) to[bend left=48] (1.5,0);
 \draw (0,0) to[bend right=48] (1.5,0);
 \fill (0,0) circle (2pt);\fill (1.5,0) circle (2pt);
 \node at (0.75,-1){1PI bubble};
 \draw (4,-0.6)--(4.8,0)--(4,0.6);
 \draw (4.8,0)--(5.6,-0.6);\draw (4.8,0)--(6.8,0.6);
 \draw (5.899,0.695) circle (0.35);\fill (6,0.36) circle (2pt);
 \fill (4.8,0) circle (2pt);\node at (5.4,-1){external-leg correction};
\end{tikzpicture}
\end{center}
At order \(\lambda^2\), the four-point function contains both the three
1PI bubbles and the four choices for a tadpole on an external leg of
the tree vertex. Merely taking a logarithm removes neither type.
Using full external propagators absorbs these leg corrections into
those propagators; it leaves the amputated interaction kernel to be
computed separately.

The \(\mathbb Z_2\) symmetry forbids a three-point vertex in this vacuum.
As a result the fully amputated connected four-point kernel is 1PI.
This special simplification does not generalize to the six-point
function: two quartic vertices joined by one line already give a
connected, reducible six-point tree.

\module{Momentum-space rules and independent loop momenta}
Fourier transform with \(\phi(x)=\intp e^{-\ii p\cdot x}\widetilde\phi(p)\).
The integral over a vertex position gives
\begin{equation}
 \int\dd^4z\,e^{-\ii(p_1+p_2+p_3+p_4)\cdot z}
 =(2\pi)^4\delta^{(4)}(p_1+p_2+p_3+p_4).
\end{equation}
Use all momenta entering each vertex, reversing a sign when an oriented
internal momentum leaves it. Lines represent a real scalar, so these
arrows track momentum, not a conserved particle charge.

\begin{keyidea}[Momentum-space Feynman rules]
\begin{enumerate}
\item An internal line of momentum \(k\) contributes
\(\ii/(k^2-m^2+\ii0)\).
\item A quartic vertex contributes \(-\ii\lambda\), with momentum
conservation at that vertex.
\item Integrate each independent loop momentum with
\(\int\dd^4\ell/(2\pi)^4\), and multiply the graph by \(1/S\).
\item Factor out the remaining overall
\((2\pi)^4\delta^{(4)}(\sum p_a)\). Include external propagators for a
correlator; omit them only when specifying an amputated quantity.
\end{enumerate}
\end{keyidea}
One may instead assign an integral to every internal line and keep all
vertex delta functions. Both procedures agree: for a connected graph
with \(I\) internal lines and \(V\) vertices, \(V-1\) independent
constraints eliminate momenta, leaving
\begin{equation}
 L=I-V+1.                                         \label{eq:loops}
\end{equation}
Thus it is the number of independent loops, rather than the number of
internal lines, that counts the integrals remaining after conservation.
For a quartic graph with \(E\) external legs,
\begin{equation}
 4V=2I+E,\qquad L=V-\frac E2+1.                  \label{eq:topology}
\end{equation}
A tadpole has \((V,I,E,L)=(1,1,2,1)\); a four-point bubble has
\((2,2,4,1)\). Both are one-loop graphs, although their powers of
\(\lambda\) differ. A six-point tree has \((2,1,6,0)\).

For example, the two-point self-energy insertion is conventionally
written \(-\ii\Sigma\). The tadpole rules give
\begin{equation}
 -\ii\Sigma_{\rm tad}(p^2)
 =\frac{-\ii\lambda}{2}\int\frac{\dd^4k}{(2\pi)^4}
 \frac{\ii}{k^2-m^2+\ii0}.                         \label{eq:sigma}
\end{equation}
No external momentum flows through its loop, so it is independent of
\(p\). The integral is ultraviolet divergent. The rules specify the
integrand and its weight; a regulator is still needed to evaluate it.

\module{Amputation and the first scattering amplitudes}
For a connected four-point function, factor out its momentum delta
function and call the remaining transform \(\widetilde G_{4,c}^{\rm red}\).
At tree level \eqref{eq:tree4} becomes
\begin{equation}
 \widetilde G_{4,c}^{\rm red}(p_1,p_2,p_3,p_4)
 =\left[\prod_{a=1}^4\frac{\ii}{p_a^2-m^2+\ii0}\right](-\ii\lambda).
 \label{eq:amputate}
\end{equation}
Amputation removes the four external propagators, leaving
\(-\ii\lambda\). It does not remove an internal propagator or impose
on-shell kinematics by itself.

To connect with particles, use relativistically normalized states
\(\langle\bm p|\bm q\rangle=2E_{\bm p}(2\pi)^3
\delta^{(3)}(\bm p-\bm q)\), and define
\begin{equation}
 \langle f|S-1|i\rangle
 =\ii(2\pi)^4\delta^{(4)}(P_f-P_i)\,\mathcal M_{fi}.
 \label{eq:Smatrix}
\end{equation}
For \(p_1+p_2\to p_3+p_4\), the all-incoming vertex momenta are
\(p_1,p_2,-p_3,-p_4\). At tree level
\begin{equation}
 \mathcal M_{2\to2}^{\rm tree}=-\lambda,\qquad
 s=(p_1+p_2)^2,\quad t=(p_1-p_3)^2,\quad u=(p_1-p_4)^2,
 \quad s+t+u=4m^2.
\end{equation}
The last identity uses equal on-shell masses. An off-shell vertex
function has no such kinematic constraint.

Beyond tree level, write the pole of the full propagator as
\begin{equation}
 \widetilde G_2(p)=\frac{\ii Z_{\rm pole}}{p^2-M^2+\ii0}
 +\text{terms regular at }p^2=M^2.                 \label{eq:pole}
\end{equation}
Here \(M\) is the physical mass. The residue \(Z_{\rm pole}\) is a
number, distinct from the functional \(Z[J]\). The LSZ prescription
removes the external poles, divides by \(\sqrt{Z_{\rm pole}}\) for each
external field, and takes the external momenta on shell. Equivalently,
if amputation is performed using full propagators and leaves the
kernel \(\ii\mathcal V_4\), then
\begin{equation}
 \mathcal M_{2\to2}=Z_{\rm pole}^{,2}
 \mathcal V_4\big|_{p_a^2=M^2}.                    \label{eq:LSZbrief}
\end{equation}
We use this pole statement here; a full derivation of LSZ can follow
later. In quartic theory \(Z_{\rm pole}=1+O(\lambda^2)\), so residue
corrections first affect this scattering amplitude at order \(\lambda^3\).
External-leg self-energies still determine the physical mass and must
not be counted again as vertex corrections.

\module{What the first loops teach us to expect}
The one-loop four-point kernel contains the three partitions found in
Section 4. For the channel with momentum \(q\) through the bubble,
\begin{equation}
 \ii\mathcal V_{4,q}^{(1)}=
 \frac{(-\ii\lambda)^2}{2}\int\frac{\dd^4k}{(2\pi)^4}
 \frac{\ii}{k^2-m^2+\ii0}
 \frac{\ii}{(k+q)^2-m^2+\ii0},\qquad q^2=s,t,u.
 \label{eq:bubble}
\end{equation}
The external propagators are absent. The \(1/2\) remains even when
the external particles are labeled. Identical-particle factors in
cross sections or phase space have a different origin.

Power counting predicts which integrals require attention. A scalar
propagator falls as \(k^{-2}\), so simultaneous scaling of all loop
momenta gives the superficial degree of divergence
\begin{equation}
 \omega=dL-2I=(d-4)V+d-\frac{d-2}{2}E.
\end{equation}
In four dimensions, \(\omega=4-E\). Two-point graphs can diverge
quadratically; four-point graphs can diverge logarithmically. Graphs
with more external legs have negative superficial degree, although
at higher loops divergent subgraphs can still occur. The tadpole and
bubble realize precisely the first two cases.

A dimensional check gives the same organizing principle:
\begin{equation}
 [\phi]=\frac{d-2}{2},\qquad [m]=1,\qquad[\lambda]=4-d.
\end{equation}
Here brackets denote mass dimension. In \(d=4\), the operators
\(1,\phi^2,(\partial\phi)^2,\phi^4\) are the local terms compatible
with the symmetry that we will need to discuss for ultraviolet
renormalization. The vacuum term drops out of normalized correlators;
the wavefunction term first becomes necessary at two loops in this model.

\begin{keyidea}[A practical order of operations]
Specify the quantity: full, connected, 1PI, or scattering. Fix its
external labels and perturbative order. Enumerate the graphs and
count their weights. Assign momenta and eliminate redundant delta
functions. Only then regulate and evaluate the remaining integrals.
Check the result against dimensions, symmetry, and the
Schwinger--Dyson equations.
\end{keyidea}
Lecture 22 will turn the divergent expressions into predictions by
fixing a physical mass and a coupling convention. This requires both
regularization, which defines an intermediate integral, and
renormalization, which relates its parameters to specified data.

\smallskip{\footnotesize\textbf{Reading.} Based on the supplied interaction
slides, with the derivations organized alongside R.~G.~Leigh,
\emph{Quantum Field Theory}, Sec.~6.5 (printed pp.~115--123), and
D.~Harlow, \emph{Quantum Field Theory I} (2024), Secs.~7 and 10.
Their metric and propagator conventions have been translated to those
of Lectures 15--20.}
\end{document}
